\label{kap:vice-promennych}
\mathquote{Příroda se směje obtížím integrace.}{Pierre-Simon Laplace}
\tikzset{pohled3d/.style={x={(-0.55cm,-0.45cm)}, y={(1cm,0cm)}, z={(0cm,1cm)}}}

Všechny funkce, které jsme dosud derivovali a integrovali, měly jednu proměnnou. Svět ale tak jednoduchý není. Teplota v místnosti závisí na třech souřadnicích a na čase. Cena bytu na rozloze, patře, poloze a stáří. Chyba neuronové sítě na milionech parametrů. V této kapitole přeneseme derivace a integrály do více rozměrů. Uvidíme, že většina myšlenek zůstane stejná, jen místo křivek budeme mít plochy, místo tečen tečné roviny a místo `sklonu' vektor, kterému se říká gradient.

\section{Funkce dvou proměnných jako krajina}

\begin{definitionbox}
  \textbf{Funkce dvou proměnných} $f\colon D\to\mathbb{R}$, $D\subseteq\mathbb{R}^2$, přiřazuje každému bodu $[x, y]$ číslo $f(x, y)$. Jejím \textbf{grafem} je plocha $z = f(x, y)$ v prostoru. \textbf{Vrstevnice} na úrovni $c$ je množina bodů roviny, kde $f(x, y) = c$.
\end{definitionbox}

Nejlepší představa: $f(x, y)$ je nadmořská výška krajiny v bodě se zeměpisnými souřadnicemi $[x, y]$. Graf je samotná krajina, vrstevnice jsou čáry, které znáte z turistických map. Kde jsou vrstevnice hustě u sebe, je svah prudký.

\begin{center}
  \begin{tikzpicture}[pohled3d, scale=1.25]
    \draw[->] (0,0,0) -- (2.4,0,0) node[below left] {\small $x$};
    \draw[->] (0,0,0) -- (0,2.4,0) node[right] {\small $y$};
    \foreach \r in {0.7,1.0,1.4} {
      \draw[colorofwisdom!60, dashed] plot[variable=\t, domain=0:360, samples=60] ({\r*cos(\t)}, {\r*sin(\t)}, 0);
      \draw[gray!60, dotted] ({\r},0,0) -- ({\r},0,{\r*\r/2});
    }
    \foreach \a in {0,30,...,330} {
      \draw[gray!70] plot[variable=\r, domain=0:1.5, samples=15] ({\r*cos(\a)}, {\r*sin(\a)}, {\r*\r/2});
    }
    \foreach \r in {0.7,1.0,1.4} {
      \draw[thick, colorofwisdom] plot[variable=\t, domain=0:360, samples=60] ({\r*cos(\t)}, {\r*sin(\t)}, {\r*\r/2});
    }
    \draw[thick] plot[variable=\t, domain=0:360, samples=60] ({1.5*cos(\t)}, {1.5*sin(\t)}, 1.125);
    \draw[->] (0,0,0) -- (0,0,1.9) node[above] {\small $z$};
    \node[font=\small] at (0,2.6,1.3) {$z = \frac{x^2 + y^2}{2}$};
    \node[font=\scriptsize, colorofwisdom] at (0,2.3,-0.3) {vrstevnice};
  \end{tikzpicture}
  \hspace{1cm}
  \begin{tikzpicture}[scale=1.15]
    \draw[->] (-1.8,0) -- (1.9,0) node[right] {\small $x$};
    \draw[->] (0,-1.8) -- (0,1.9) node[above] {\small $y$};
    \foreach \r/\c in {0.7/0{,}25, 1.0/0{,}5, 1.4/1} {
      \draw[thick, colorofwisdom] (0,0) circle (\r);
      \node[font=\scriptsize, fill=white, inner sep=1pt] at ({\r*cos(45)},{\r*sin(45)}) {$\c$};
    }
    \node[font=\scriptsize, align=center] at (0,-2.3) {mapa: vrstevnice $x^2 + y^2 = 2c$};
  \end{tikzpicture}
\end{center}

Vlevo je graf funkce $f(x, y) = \frac{x^2 + y^2}{2}$, tedy `miska', vpravo její mapa. Vrstevnice jsou kružnice. Čím výš, tím jsou si blíž, protože miska je dál od středu strmější.

\begin{center}
  \begin{tikzpicture}[pohled3d, scale=1.3]
    \foreach \c in {-1,-0.5,...,1} {
      \draw[gray!80] plot[variable=\t, domain=-1:1, samples=20] ({\c}, {\t}, {0.8*(\c*\c-\t*\t)});
      \draw[colorofwisdom!80] plot[variable=\t, domain=-1:1, samples=20] ({\t}, {\c}, {0.8*(\t*\t-\c*\c)});
    }
    \fill (0,0,0) circle (1.5pt);
    \node[font=\small] at (0,1.9,0.9) {$z = x^2 - y^2$};
  \end{tikzpicture}
  \hspace{1.2cm}
  \begin{tikzpicture}[scale=1]
    \clip (-1.7,-1.7) rectangle (1.7,1.7);
    \foreach \c in {0.3,0.8,1.5} {
      \draw[thick, colorofwisdom, domain=-1.3:1.3, samples=40] plot ({sqrt(\c)*cosh(\x)}, {sqrt(\c)*sinh(\x)});
      \draw[thick, colorofwisdom, domain=-1.3:1.3, samples=40] plot ({-sqrt(\c)*cosh(\x)}, {sqrt(\c)*sinh(\x)});
      \draw[thick, gray, domain=-1.3:1.3, samples=40] plot ({sqrt(\c)*sinh(\x)}, {sqrt(\c)*cosh(\x)});
      \draw[thick, gray, domain=-1.3:1.3, samples=40] plot ({sqrt(\c)*sinh(\x)}, {-sqrt(\c)*cosh(\x)});
    }
    \draw[dashed] (-1.7,-1.7) -- (1.7,1.7) (-1.7,1.7) -- (1.7,-1.7);
    \draw[->] (-1.7,0) -- (1.7,0);
    \draw[->] (0,-1.7) -- (0,1.7);
  \end{tikzpicture}
\end{center}

Funkce $f(x, y) = x^2 - y^2$ je \textbf{sedlo}, tvar, který znáte z koňského sedla nebo z chipsů Pringles. Ve směru osy $x$ jde nahoru, ve směru osy $y$ dolů. Její vrstevnice jsou hyperboly z kapitoly~\ref{kap:kuzelosecky}: barevné jsou kladné úrovně (`kopce' vlevo a vpravo), šedé záporné (`údolí' nahoře a dole). Bod $[0, 0]$ není ani vrchol, ani dno, a přesto je v něm krajina `vodorovná'.

\section{Parciální derivace}

Jak změřit sklon krajiny? Záleží na tom, kterým směrem jdeme. Nejjednodušší je jít rovnoběžně s osami.

\begin{definitionbox}
  \textbf{Parciální derivace} funkce $f(x, y)$ podle $x$ v bodě $[a, b]$ je
  \begin{equation*}
    \frac{\partial f}{\partial x}(a, b) = \lim_{h\to 0}\frac{f(a + h, b) - f(a, b)}{h}.
  \end{equation*}
  Počítá se jako obyčejná derivace podle $x$, přičemž $y$ považujeme za konstantu. Podobně $\frac{\partial f}{\partial y}$. Značí se také $f_x$, $f_y$.
\end{definitionbox}

Geometricky: krajinu `rozřízneme' svislou rovinou. Pro $f_x$ rovinou $y = b$, pro $f_y$ rovinou $x = a$ (na obrázku). Řez je obyčejná křivka a parciální derivace je sklon její tečny. Jako byste šli přesně na východ (nebo na sever) a měřili, jak strmě stoupáte.

\begin{center}
  \begin{tikzpicture}[pohled3d, scale=1.5]
    \foreach \a in {0,30,...,330} {
      \draw[gray!50] plot[variable=\r, domain=0:1.5, samples=15] ({\r*cos(\a)}, {\r*sin(\a)}, {\r*\r/2});
    }
    \draw[gray!60] plot[variable=\t, domain=0:360, samples=60] ({1.5*cos(\t)}, {1.5*sin(\t)}, 1.125);
    \fill[colorofwisdom!12, opacity=0.7] (0.6,-1.7,0) -- (0.6,1.7,0) -- (0.6,1.7,1.6) -- (0.6,-1.7,1.6) -- cycle;
    \draw[colorofwisdom!50] (0.6,-1.7,0) -- (0.6,1.7,0) -- (0.6,1.7,1.6) -- (0.6,-1.7,1.6) -- cycle;
    \draw[very thick, colorofwisdom] plot[variable=\t, domain=-1.38:1.38, samples=40] (0.6, {\t}, {(\t*\t+0.36)/2});
    \draw[thick] (0.6,0.1,-0.06) -- (0.6,1.5,1.06);
    \fill (0.6,0.8,0.5) circle (1.3pt);
    \node[font=\scriptsize, anchor=south] at (0.6,-1.2,1.6) {řez $x = a$};
    \node[font=\scriptsize, anchor=west] at (0.6,1.55,1.06) {tečna se sklonem $f_y(a, b)$};
  \end{tikzpicture}
\end{center}

\begin{example}
  Pro $f(x, y) = x^2y + \sin(xy)$:
  \begin{equation*}
    \frac{\partial f}{\partial x} = 2xy + y\cos(xy),\qquad\frac{\partial f}{\partial y} = x^2 + x\cos(xy).
  \end{equation*}
  Derivujeme-li podle $x$, je $y$ jen číslo, třeba $5$. Pak $x^2\cdot 5$ má derivaci $10x$, tedy $2xy$, a $\sin(5x)$ má derivaci $5\cos(5x)$, tedy $y\cos(xy)$.
\end{example}

Derivovat můžeme i podruhé: $f_{xx}$, $f_{yy}$ a smíšené $f_{xy}$, $f_{yx}$. Pro rozumné funkce na pořadí nezáleží, $f_{xy} = f_{yx}$ (\textbf{Schwarzova věta}). V našem příkladu $f_{xy} = 2x + \cos(xy) - xy\sin(xy) = f_{yx}$. Zkontrolujte si to.

\section{Gradient: kudy nahoru}

\begin{definitionbox}
  \textbf{Gradient} funkce $f$ v bodě $[a, b]$ je vektor parciálních derivací
  \begin{equation*}
    \nabla f(a, b) = \big(f_x(a, b),\ f_y(a, b)\big).
  \end{equation*}
  \textbf{Směrová derivace} ve směru jednotkového vektoru $\vec{u}$ je
  \begin{equation*}
    D_{\vec{u}}f = \lim_{h\to 0}\frac{f([a, b] + h\vec{u}) - f(a, b)}{h} = \nabla f\cdot\vec{u}.
  \end{equation*}
\end{definitionbox}

Druhá rovnost (pro hladké funkce) říká, že sklon v libovolném směru dostaneme ze dvou čísel $f_x$ a $f_y$. Krok $h\vec{u} = (hu_1, hu_2)$ se skládá z kroku $hu_1$ na východ a $hu_2$ na sever. Výška se při něm změní přibližně o $f_x\cdot hu_1 + f_y\cdot hu_2 = h(\nabla f\cdot\vec{u})$.

Ze skalárního součinu plynou dvě klíčové vlastnosti:
\begin{itemize}
  \item $D_{\vec{u}}f = |\nabla f|\cos\theta$, kde $\theta$ je úhel mezi $\vec{u}$ a gradientem. Největší je pro $\theta = 0$. \textbf{Gradient ukazuje směr nejprudšího stoupání} a jeho délka je velikost toho stoupání. Opačný směr, $-\nabla f$, je nejprudší klesání. Kam by stekla voda.
  \item Podél vrstevnice se výška nemění, takže směrová derivace podél ní je nulová: $\nabla f\cdot\vec{t} = 0$. \textbf{Gradient je kolmý na vrstevnice.}
\end{itemize}

\begin{center}
  \begin{tikzpicture}[scale=1.3]
    \foreach \c in {0.5,1.5,3} {
      \draw[thick, colorofwisdom!70] plot[variable=\t, domain=0:360, samples=80] ({sqrt(\c)*cos(\t)}, {sqrt(\c/2)*sin(\t)});
    }
    \foreach \c in {0.5,1.5} {
      \foreach \t in {20,70,110,160,200,250,290,340} {
        \pgfmathsetmacro{\px}{sqrt(\c)*cos(\t)}
        \pgfmathsetmacro{\py}{sqrt(\c/2)*sin(\t)}
        \draw[->, thick] (\px,\py) -- ({\px+0.18*2*\px},{\py+0.18*4*\py});
      }
    }
    \draw[->, gray] (-2.2,0) -- (2.3,0) node[right] {\small $x$};
    \draw[->, gray] (0,-1.5) -- (0,1.6) node[above] {\small $y$};
    \node[font=\scriptsize, align=center] at (0,-1.95) {vrstevnice $x^2 + 2y^2 = c$ a gradienty $\nabla f = (2x, 4y)$};
  \end{tikzpicture}
\end{center}

Šipky gradientu míří kolmo ven z vrstevnic, tedy do kopce. Všimněte si, že nemíří do středu: v protáhlém údolí nevede nejprudší směr přímo ke dnu. Tohle je mimochodem základ \textbf{gradientního sestupu}, algoritmu, kterým se učí neuronové sítě. Chybu sítě si představíme jako krajinu v mnoha rozměrech a opakovaně uděláme malý krok ve směru $-\nabla f$. Krok po kroku sejdeme do údolí.

\section{Tečná rovina a lineární přiblížení}

Tečna ke křivce je přímka, která se ke křivce v jednom bodě nejvíc `přimkne'. U plochy hraje stejnou roli \textbf{tečná rovina}. Když se díváte na krajinu z dostatečné blízkosti, vypadá rovinná. Proto lidé dlouho věřili, že Země je placatá.

\begin{definitionbox}
  \textbf{Tečná rovina} ke grafu $z = f(x, y)$ v bodě $[a, b, f(a, b)]$:
  \begin{equation*}
    z = f(a, b) + f_x(a, b)(x - a) + f_y(a, b)(y - b).
  \end{equation*}
  Pro body blízko $[a, b]$ tedy platí \textbf{lineární přiblížení} (totální diferenciál)
  \begin{equation*}
    \Delta f\approx f_x\,\Delta x + f_y\,\Delta y.
  \end{equation*}
\end{definitionbox}

\begin{example}
  Válec má poloměr $r = 5$ cm a výšku $v = 10$ cm, oboje změřené s chybou nejvýš $0{,}1$ cm. Jak velká může být chyba objemu $V = \pi r^2v$? Platí $V_r = 2\pi rv = 100\pi$ a $V_v = \pi r^2 = 25\pi$, takže
  \begin{equation*}
    |\Delta V|\lessapprox 100\pi\cdot 0{,}1 + 25\pi\cdot 0{,}1 = 12{,}5\pi\approx 39\ \text{cm}^3,
  \end{equation*}
  což je asi $5\,\%$ objemu $250\pi\approx 785$ cm$^3$. Chyba v poloměru je čtyřikrát nebezpečnější než chyba ve výšce, protože poloměr je ve vzorci na druhou. Takhle fyzici počítají chyby měření.
\end{example}

\section{Extrémy}

Kde má krajina vrchol nebo dno, je tečná rovina vodorovná, tedy obě parciální derivace jsou nulové. Bodům s $\nabla f = \vec{o}$ říkáme \textbf{stacionární}. Jenže vodorovná je i rovina v sedle. Jak je rozlišit?

V jedné proměnné rozhodovala druhá derivace, tedy jestli je graf vypouklý nahoru nebo dolů. Ve dvou proměnných máme tři druhé derivace $f_{xx}$, $f_{xy}$, $f_{yy}$ a krajina může být v různých směrech zakřivená různě. Kolem stacionárního bodu se funkce chová jako kvadratický výraz
\begin{equation*}
  f(a + h, b + k)\approx f(a, b) + \frac{1}{2}\left(f_{xx}h^2 + 2f_{xy}hk + f_{yy}k^2\right),
\end{equation*}
tedy jako miska, obrácená miska, nebo sedlo. Který z těch tvarů to je, rozhoduje znaménko jistého determinantu.

\begin{definitionbox}
  Nechť $\nabla f(a, b) = \vec{o}$ a označme $D = f_{xx}f_{yy} - f_{xy}^2$ (vše v bodě $[a, b]$).
  \begin{itemize}
    \item $D > 0$ a $f_{xx} > 0$: \textbf{lokální minimum} (miska),
    \item $D > 0$ a $f_{xx} < 0$: \textbf{lokální maximum} (obrácená miska),
    \item $D < 0$: \textbf{sedlový bod},
    \item $D = 0$: test nerozhodne.
  \end{itemize}
\end{definitionbox}

Číslo $D$ je determinant symetrické matice $\begin{pmatrix} f_{xx} & f_{xy}\\ f_{xy} & f_{yy}\end{pmatrix}$, které se říká \textbf{Hessova matice}. Podle spektrální věty z kapitoly~\ref{kap:linearni-algebra} má dvě reálná vlastní čísla, která udávají zakřivení ve dvou kolmých `hlavních' směrech, a determinant je jejich součin. Kladný součin znamená stejná znaménka (obě nahoru nebo obě dolů), záporný různá znaménka, tedy sedlo.

\begin{example}
  Najděme extrémy funkce $f(x, y) = x^3 - 3x + y^2$. Stacionární body: $f_x = 3x^2 - 3 = 0$ a $f_y = 2y = 0$, tedy $[1, 0]$ a $[-1, 0]$. Druhé derivace: $f_{xx} = 6x$, $f_{yy} = 2$, $f_{xy} = 0$.
  \begin{itemize}
    \item V $[1, 0]$: $D = 6\cdot 2 = 12 > 0$, $f_{xx} = 6 > 0$, lokální minimum $f(1, 0) = -2$.
    \item V $[-1, 0]$: $D = -12 < 0$, sedlo.
  \end{itemize}
  Globální minimum funkce nemá, protože pro $x\to-\infty$ jde $x^3$ do $-\infty$.
\end{example}

\subsection{Vázané extrémy a Lagrangeovy multiplikátory}

Často hledáme extrém jen na nějaké křivce: maximum zisku při daném rozpočtu, nejkratší cestu po povrchu Země. Úloha zní: najdi extrém $f(x, y)$ za podmínky $g(x, y) = 0$.

Představte si, že jdete po turistické stezce (to je podmínka $g = 0$) a hledáte na ní nejvyšší místo. Dokud stezka vrstevnice protíná, můžete ještě vystoupat výš. V nejvyšším bodě se stezka vrstevnice jen \textit{dotkne}. Tam je stezka rovnoběžná s vrstevnicí, a tedy jsou rovnoběžné i jejich normály, gradienty $\nabla f$ a $\nabla g$.

\begin{center}
  \begin{tikzpicture}[scale=1.6]
    \foreach \c in {-1,-0.5,0,0.5,1,1.4142,2} {
      \draw[colorofwisdom!60] ({\c/2-1.1},{\c/2+1.1}) -- ({\c/2+1.1},{\c/2-1.1});
    }
    \node[font=\scriptsize, colorofwisdom] at (1.85,0.05) {$x + y = c$};
    \draw[very thick] (0,0) circle (1);
    \node[font=\scriptsize] at (-0.95,-0.95) {$x^2 + y^2 = 1$};
    \fill (0.7071,0.7071) circle (1.3pt);
    \draw[->, thick, colorofwisdom] (0.7071,0.7071) -- (1.1071,1.1071) node[above right, font=\scriptsize] {$\nabla f\parallel\nabla g$};
    \draw[->, gray] (-1.5,0) -- (1.6,0) node[right] {\small $x$};
    \draw[->, gray] (0,-1.5) -- (0,1.6) node[above] {\small $y$};
  \end{tikzpicture}
\end{center}

\begin{definitionbox}
  \textbf{Metoda Lagrangeových multiplikátorů}: V bodě vázaného extrému funkce $f$ za podmínky $g = 0$ (kde $\nabla g\neq\vec{o}$) existuje číslo $\lambda$ takové, že
  \begin{equation*}
    \nabla f = \lambda\nabla g,\qquad g = 0.
  \end{equation*}
\end{definitionbox}

\begin{example}
  Maximum $f(x, y) = x + y$ na kružnici $x^2 + y^2 = 1$ (obrázek nahoře). Podmínka $(1, 1) = \lambda(2x, 2y)$ dává $x = y = \frac{1}{2\lambda}$, dosazením do kružnice $x = y = \pm\frac{\sqrt{2}}{2}$. Maximum je $\sqrt{2}$ v bodě $\left[\frac{\sqrt{2}}{2}, \frac{\sqrt{2}}{2}\right]$, minimum $-\sqrt{2}$ v bodě opačném. Na obrázku je vidět, že přímka $x + y = \sqrt{2}$ se kružnice právě dotýká.
\end{example}

\section{Dvojný integrál}

Integrál funkce jedné proměnné byl plocha pod grafem. Integrál funkce dvou proměnných přes oblast $M$ v rovině je \textbf{objem pod grafem}. Definice je stejná: rozdělíme $M$ na malé obdélníčky o obsahu $\Delta S$, nad každým postavíme hranol s výškou $f$ a objemy sečteme. V limitě dostaneme $\iint_M f(x, y)\,\mathrm{d}S$.

Jak se takový integrál počítá? Úplně stejně, jako se krájí chleba.

\begin{definitionbox}
  \textbf{Fubiniho věta}: Pro spojitou funkci na obdélníku $[a, b]\times[c, d]$ je
  \begin{equation*}
    \iint_{[a,b]\times[c,d]}f(x, y)\,\mathrm{d}S = \int_a^b\left(\int_c^d f(x, y)\,\mathrm{d}y\right)\mathrm{d}x = \int_c^d\left(\int_a^b f(x, y)\,\mathrm{d}x\right)\mathrm{d}y.
  \end{equation*}
\end{definitionbox}

Vnitřní integrál $\int_c^d f(x, y)\,\mathrm{d}y$ je obsah jednoho `krajíce', řezu tělesa rovinou s pevným $x$. Vnější integrál krajíce sečte. Je jedno, jestli krájíme podél nebo napříč, chleba má pořád stejný objem. Je to tatáž myšlenka jako Cavalieriho princip ze stereometrie.

\begin{center}
  \begin{tikzpicture}[pohled3d, scale=1.4]
    \draw[->] (0,0,0) -- (2.6,0,0) node[below left] {\small $x$};
    \draw[->] (0,0,0) -- (0,2.8,0) node[right] {\small $y$};
    \draw[->] (0,0,0) -- (0,0,2.9) node[above] {\small $z$};
    \draw[gray] (2,0,0) -- (2,2,0) -- (0,2,0);
    \foreach \x in {0,0.5,1,1.5,2} {
      \draw[gray!70] plot[variable=\t, domain=0:2, samples=10] ({\x}, {\t}, {0.6+0.5*\x*\t});
    }
    \foreach \y in {0,0.5,1,1.5,2} {
      \draw[gray!70] plot[variable=\t, domain=0:2, samples=10] ({\t}, {\y}, {0.6+0.5*\t*\y});
    }
    \fill[colorofwisdom!45, opacity=0.85] (1.2,0,0) -- (1.2,2,0) -- plot[variable=\t, domain=2:0, samples=10] (1.2, {\t}, {0.6+0.6*\t}) -- cycle;
    \draw[colorofwisdom, thick] (1.2,0,0) -- (1.2,2,0) -- (1.2,2,1.8) -- (1.2,0,0.6) -- cycle;
    \foreach \p in {(2,0,0),(2,2,0),(0,2,0)} \draw[gray, dotted] \p -- ++(0,0,0.6);
    \node[font=\scriptsize, anchor=west] at (1.2,2.1,1.7) {krajíc $\int f(x_0, y)\,\mathrm{d}y$};
  \end{tikzpicture}
\end{center}

\begin{example}
  $\displaystyle\iint_{[0,1]\times[0,2]}xy\,\mathrm{d}S = \int_0^1\left(\int_0^2 xy\,\mathrm{d}y\right)\mathrm{d}x = \int_0^1 x\left[\frac{y^2}{2}\right]_0^2\mathrm{d}x = \int_0^1 2x\,\mathrm{d}x = 1$.
\end{example}

\subsection{Polární souřadnice}

Pro kruhové oblasti se hodí popisovat body vzdáleností $r$ od počátku a úhlem $\varphi$: $x = r\cos\varphi$, $y = r\sin\varphi$. Kousek oblasti mezi $r$ a $r + \mathrm{d}r$ a mezi $\varphi$ a $\varphi + \mathrm{d}\varphi$ je skoro obdélník se stranami $\mathrm{d}r$ a $r\,\mathrm{d}\varphi$ (délka oblouku je poloměr krát úhel). Jeho obsah je tedy
\begin{equation*}
  \mathrm{d}S = r\,\mathrm{d}r\,\mathrm{d}\varphi.
\end{equation*}

\begin{center}
  \begin{tikzpicture}[scale=1.6]
    \draw[->, gray] (-0.2,0) -- (2.6,0);
    \draw[->, gray] (0,-0.2) -- (0,2.3);
    \foreach \r in {0.5,1,1.5,2} \draw[gray!50] (\r,0) arc (0:90:\r);
    \foreach \a in {15,30,...,75} \draw[gray!50] (0,0) -- (\a:2.1);
    \fill[colorofwisdom!45] (30:1.5) arc (30:45:1.5) -- (45:2) arc (45:30:2) -- cycle;
    \draw[thick, colorofwisdom] (30:1.5) arc (30:45:1.5) -- (45:2) arc (45:30:2) -- cycle;
    \draw[decorate, decoration={brace, amplitude=3pt}] (30:1.5) -- (30:2) node[midway, below right, font=\scriptsize] {$\mathrm{d}r$};
    \node[font=\scriptsize] at (37.5:2.35) {$r\,\mathrm{d}\varphi$};
    \draw (0.5,0) arc (0:30:0.5);
    \node[font=\scriptsize] at (15:0.7) {$\varphi$};
  \end{tikzpicture}
\end{center}

Ten faktor $r$ je důležitý: dílky dál od středu jsou větší. Vyplatí se to hned v jednom z nejkrásnějších výpočtů celé analýzy.

\section{Odbočka --- Gaussův integrál}

Funkce $\mathrm{e}^{-x^2}$ je základem normálního rozdělení z kapitoly~\ref{kap:pravdepodobnost}. Její primitivní funkci nelze zapsat pomocí elementárních funkcí (zmínili jsme to u základní věty integrálního počtu). Přesto spočítáme
\begin{equation*}
  I = \int_{-\infty}^{\infty}\mathrm{e}^{-x^2}\,\mathrm{d}x.
\end{equation*}

\begin{center}
  \begin{tikzpicture}[x=1.3cm, y=2.4cm]
    \fill[colorofwisdom!25, domain=-2.6:2.6, samples=80] (-2.6,0) -- plot (\x, {exp(-\x*\x)}) -- (2.6,0) -- cycle;
    \draw[very thick, colorofwisdom, domain=-2.8:2.8, samples=100] plot (\x, {exp(-\x*\x)});
    \draw[->] (-3,0) -- (3.1,0) node[right] {\small $x$};
    \draw[->] (0,0) -- (0,1.25) node[above] {\small $y$};
    \node[font=\scriptsize, left] at (0,1) {$1$};
    \node[font=\small] at (0,0.35) {obsah $= \sqrt{\pi}\approx 1{,}772$};
    \node[font=\small, colorofwisdom, anchor=west] at (0.8,0.75) {$y = \mathrm{e}^{-x^2}$};
  \end{tikzpicture}
\end{center}

Trik je neuvěřitelně drzý: místo $I$ spočítáme $I^2$ a podíváme se na něj jako na dvojný integrál přes celou rovinu.
\begin{equation*}
  I^2 = \int_{-\infty}^{\infty}\mathrm{e}^{-x^2}\,\mathrm{d}x\cdot\int_{-\infty}^{\infty}\mathrm{e}^{-y^2}\,\mathrm{d}y = \iint_{\mathbb{R}^2}\mathrm{e}^{-(x^2 + y^2)}\,\mathrm{d}S.
\end{equation*}
Funkce $\mathrm{e}^{-(x^2 + y^2)}$ je `zvon', který závisí jen na vzdálenosti $r$ od počátku, protože $x^2 + y^2 = r^2$. Je souměrný kolem osy $z$, a tak se vyplatí krájet ho kolem dokola.

\begin{center}
  \begin{tikzpicture}[pohled3d, scale=1.35]
    \draw[->] (0,0,0) -- (2.8,0,0) node[below left] {\small $x$};
    \draw[->] (0,0,0) -- (0,2.8,0) node[right] {\small $y$};
    \foreach \r in {0.4,1.2,1.6,2.0} {
      \draw[gray!45] plot[variable=\t, domain=0:360, samples=60] ({\r*cos(\t)}, {\r*sin(\t)}, {2*exp(-\r*\r)});
    }
    \foreach \a in {0,30,...,330} {
      \draw[gray!60] plot[variable=\r, domain=0:2.2, samples=30] ({\r*cos(\a)}, {\r*sin(\a)}, {2*exp(-\r*\r)});
    }
    \fill[colorofwisdom!30, opacity=0.6] plot[variable=\t, domain=0:360, samples=60] ({0.8*cos(\t)}, {0.8*sin(\t)}, 0);
    \draw[colorofwisdom, thick] plot[variable=\t, domain=0:360, samples=60] ({0.8*cos(\t)}, {0.8*sin(\t)}, 0);
    \draw[colorofwisdom, very thick] plot[variable=\t, domain=0:360, samples=60] ({0.8*cos(\t)}, {0.8*sin(\t)}, 1.0546);
    \draw[colorofwisdom, thick] ({0.8*cos(118.8)}, {0.8*sin(118.8)}, 0) -- ({0.8*cos(118.8)}, {0.8*sin(118.8)}, 1.0546);
    \draw[colorofwisdom, thick] ({0.8*cos(-61.2)}, {0.8*sin(-61.2)}, 0) -- ({0.8*cos(-61.2)}, {0.8*sin(-61.2)}, 1.0546);
    \draw[->] (0,0,0) -- (0,0,2.5) node[above] {\small $z$};
    \node[font=\small, anchor=west] at (0,2.3,1.7) {$z = \mathrm{e}^{-(x^2 + y^2)} = \mathrm{e}^{-r^2}$};
    \node[font=\scriptsize, colorofwisdom, anchor=west] at ({0.8*cos(-61.2)}, {0.8*sin(-61.2)}, 0.5) {slupka o poloměru $r$};
  \end{tikzpicture}
  \hspace{0.6cm}
  \begin{tikzpicture}[scale=1.1]
    \foreach \r/\k in {2.0/4, 1.8/6, 1.6/10, 1.4/15, 1.2/22, 1.0/30, 0.8/40, 0.6/50, 0.4/58, 0.2/63} {
      \fill[colorofwisdom!\k] (0,0) circle (\r);
    }
    \fill[colorofwisdom!90!black, even odd rule] (0,0) circle (1.0) (0,0) circle (0.8);
    \draw[->, thick] (0,0) -- (-30:0.8) node[midway, below left, font=\scriptsize] {$r$};
    \draw[decorate, decoration={brace, amplitude=2pt}] (40:1.02) -- (40:0.78);
    \node[font=\scriptsize] at (40:1.3) {$\mathrm{d}r$};
    \node[font=\scriptsize, align=center] at (0,-2.45) {pohled shora: mezikruží\\o obsahu $\approx 2\pi r\,\mathrm{d}r$};
  \end{tikzpicture}
\end{center}

Představte si zvon rozkrájený na tenké válcové slupky, jako cibuli. Slupka o poloměru $r$ a tloušťce $\mathrm{d}r$ je vysoká $\mathrm{e}^{-r^2}$. Když ji rozřízneme a rozvineme, dostaneme skoro kvádr o rozměrech $2\pi r\times\mathrm{d}r\times\mathrm{e}^{-r^2}$. Celkový objem je součet všech slupek:
\begin{equation*}
  I^2 = \int_0^{\infty}\underbrace{2\pi r}_{\text{obvod}}\cdot\underbrace{\mathrm{e}^{-r^2}}_{\text{výška}}\,\mathrm{d}r.
\end{equation*}
Faktor $r$ tedy není žádný formální trik, je to obvod slupky. Vnější slupky jsou delší, a proto do objemu přispívají víc.

Totéž v polárních souřadnicích:
\begin{equation*}
  I^2 = \int_0^{2\pi}\int_0^{\infty}\mathrm{e}^{-r^2}\,r\,\mathrm{d}r\,\mathrm{d}\varphi = 2\pi\left[-\frac{1}{2}\mathrm{e}^{-r^2}\right]_0^{\infty} = 2\pi\cdot\frac{1}{2} = \pi.
\end{equation*}
Faktor $r$ z polárních souřadnic je přesně ta derivace vnitřní funkce, kterou potřebujeme k substituci. Bez něj by to nešlo. Tedy
\begin{equation*}
  \int_{-\infty}^{\infty}\mathrm{e}^{-x^2}\,\mathrm{d}x = \sqrt{\pi}.
\end{equation*}
Integrál, který v jedné dimenzi spočítat nejde, se ve dvou dimenzích spočítá na tři řádky. Proto má hustota normálního rozdělení ve jmenovateli $\sqrt{2\pi}$. A proto je $\Gamma\!\left(\frac{1}{2}\right) = \sqrt{\pi}$, jak jsme viděli v odbočce o faktoriálu v kapitole~\ref{kap:kombinatorika}.

\section{Vektorová pole}

Na závěr ještě jeden pojem, bez kterého se fyzika neobejde. \textbf{Vektorové pole} přiřazuje každému bodu roviny (nebo prostoru) vektor. Rychlost proudící vody, síla větru na meteorologické mapě, gravitační nebo elektrické pole. Gradient funkce je vektorové pole, ale ne každé pole je gradient.

\begin{center}
  \begin{tikzpicture}[scale=0.85]
    \foreach \x in {-2,-1,...,2} \foreach \y in {-2,-1,...,2} {
      \draw[->, colorofwisdom, thick] (\x,\y) -- ({\x+0.18*\x},{\y+0.18*\y});
    }
    \node[font=\small] at (0,-2.8) {$\vec{F} = (x, y)$: zdroj};
    \begin{scope}[xshift=6cm]
      \foreach \x in {-2,-1,...,2} \foreach \y in {-2,-1,...,2} {
        \draw[->, colorofwisdom, thick] (\x,\y) -- ({\x-0.18*\y},{\y+0.18*\x});
      }
      \node[font=\small] at (0,-2.8) {$\vec{F} = (-y, x)$: vír};
    \end{scope}
  \end{tikzpicture}
\end{center}

Pole vlevo `vyvěrá' z počátku jako voda z pramene. Jeho \textbf{divergence} $\frac{\partial F_1}{\partial x} + \frac{\partial F_2}{\partial y} = 2$ je kladná: v každém bodě voda přibývá. Pole vpravo se točí. Jeho divergence je nulová, ale jeho \textbf{rotace} $\frac{\partial F_2}{\partial x} - \frac{\partial F_1}{\partial y} = 2$ ne. Kdybyste do něj položili malé mlýnské kolečko, roztočilo by se. Pole vpravo není gradientem žádné funkce. Kdyby bylo, šli bychom po kružnici pořád `do kopce' a vrátili se do výchozího bodu výš, než jsme začali, což je nesmysl. Divergence a rotace jsou jazyk, kterým jsou napsány Maxwellovy rovnice elektromagnetismu.

\begin{problem}
  \begin{enumerate}[label=(\alph*)]
    \item Spočítejte parciální derivace funkce $f(x, y) = \mathrm{e}^{xy}\ln y$.
    \item Ve kterém směru roste funkce $f(x, y) = x^2y$ v bodě $[1, 2]$ nejrychleji a jak rychle?
    \item Najděte a klasifikujte stacionární body funkce $f(x, y) = x^2 + xy + y^2 - 3x$.
    \item Pomocí Lagrangeových multiplikátorů najděte obdélník s obvodem $20$ a s největším obsahem.
    \item Spočítejte objem pod grafem $z = x + y$ nad čtvercem $[0, 1]\times[0, 1]$.
  \end{enumerate}
\end{problem}

\begin{solution}
  \begin{enumerate}[label=(\alph*)]
    \item $f_x = y\,\mathrm{e}^{xy}\ln y$, $f_y = x\,\mathrm{e}^{xy}\ln y + \frac{\mathrm{e}^{xy}}{y}$.
    \item $\nabla f = (2xy, x^2) = (4, 1)$. Nejrychleji ve směru $(4, 1)$, rychlostí $|\nabla f| = \sqrt{17}$.
    \item $f_x = 2x + y - 3 = 0$, $f_y = x + 2y = 0$, tedy $[2, -1]$. $D = 2\cdot 2 - 1 = 3 > 0$, $f_{xx} = 2 > 0$, minimum $f(2, -1) = -3$.
    \item Maximalizujeme $xy$ za podmínky $2x + 2y = 20$. $(y, x) = \lambda(2, 2)$ dává $x = y = 5$. Čtverec s obsahem $25$.
    \item $\int_0^1\int_0^1(x + y)\,\mathrm{d}y\,\mathrm{d}x = \int_0^1\left(x + \frac{1}{2}\right)\mathrm{d}x = 1$.
  \end{enumerate}
\end{solution}

\section{Příběh --- Hlava, která se znovu nenarodí}

Joseph-Louis Lagrange se narodil v roce 1736 v Turíně jako Giuseppe Lodovico Lagrangia. V devatenácti letech poslal Eulerovi dopis s novou metodou, ze které vznikl variační počet, a Euler, tehdy nejslavnější matematik světa, se podle legendy se zveřejněním vlastní práce na stejné téma zdržel, aby mladík dostal zásluhy. Později Lagrange nahradil Eulera v Berlíně a v roce 1787 přesídlil do Paříže. Dva roky před revolucí.

Za jakobínského teroru byl přijat dekret o vypovězení cizinců. Lagrange, rodilý Ital, byl z něj výslovně vyjmut, a místo do vyhnanství šel pracovat do komise, která připravovala metrickou soustavu. Jeho přítel, chemik Antoine Lavoisier, takové štěstí neměl. Lavoisier byl kromě chemie i výběrčím daní, a v roce 1794 ho za to revoluční tribunál poslal pod gilotinu. Lagrange druhý den řekl: `Stačil jim okamžik, aby tu hlavu sťali, a sto let možná nestačí, aby se zrodila jiná, jako byla ona.'

Lagrange revoluci přežil, Napoleon z něj udělal hraběte a senátora a zemřel v roce 1813 v Paříži. Na vzniku metru a kilogramu, které dnes používá skoro celý svět, má tak nemalý podíl.

\section{Shrnutí a otázky}

\begin{problem}
  Odpovězte vlastními slovy:
  \begin{enumerate}
    \item Jak si představit funkci dvou proměnných? Co jsou vrstevnice?
    \item Co je parciální derivace a jak ji vidíme na grafu?
    \item Co je gradient? Proč ukazuje směr nejprudšího stoupání a proč je kolmý na vrstevnice?
    \item Co je tečná rovina a k čemu se hodí lineární přiblížení?
    \item Jak se hledají a klasifikují lokální extrémy? Co je sedlový bod?
    \item Vysvětlete geometricky metodu Lagrangeových multiplikátorů.
    \item Co říká Fubiniho věta? Proč je v polárních souřadnicích $\mathrm{d}S = r\,\mathrm{d}r\,\mathrm{d}\varphi$?
    \item Spočítejte Gaussův integrál.
    \item Co je vektorové pole, divergence a rotace?
  \end{enumerate}
\end{problem}
